\vspace{0.5em}\noindent\textbf{Sự kiện ``AWS: Enterprise Cloud Architectures and Industry
Application featuring Cloud Kinetics \& Renova
Cloud''}\par

\vspace{0.5em}\noindent\textbf{Mục Đích Của Sự
Kiện}\par

\begin{itemize}
\tightlist
\item
  Giới thiệu thị trường việc làm ngành IT hiện tại và xu hướng việc làm
  tương lai
\item
  Tìm hiểu về sự khác biệt trong về yêu cầu, môi trường làm việc giữa
  nhà trường và doanh nghiệp
\item
  Đưa ra một số lời khuyên về kỹ năng mềm, đặc biệt là kỹ năng giao tiếp
  hiệu quả trong môi trường công sở
\item
  Những kỹ năng cần thiết trong thời đại AI
\end{itemize}

\vspace{0.5em}\noindent\textbf{Danh Sách Diễn
Giả}\par

\begin{itemize}
\tightlist
\item
  \textbf{Nguyen Gia Hung} - Head of Solution Architect in Vietnam \&
  Cambodia, AWS
\item
  \textbf{Khang Nguyen} - Solutions Architect, Cloud Kinetics
\item
  \textbf{Nhu Tran} - Account Manager, AWS
\item
  \textbf{Vinh Banh} - Senior Data Engineer, Renova Cloud \#\#\# Nội
  Dung Nổi Bật
\end{itemize}

\vspace{0.5em}\noindent\textbf{Giới thiệu về Thị trường việc làm IT hiện nay và xu hướng
nghề nghiệp tương
lai}\par

\begin{itemize}
\tightlist
\item
  Xu hướng AI đang thay đổi cách chúng ta làm việc, yêu cầu kỹ năng mới
  và phương pháp tiếp cận mới.
\item
  Sự dịch chuyển từ on-premise sang cloud-native và hybrid cloud.
\item
  Sự chuyển dịch từ mô hình công việc theo vai trò cố định (role-based
  jobs) sang mô hình công việc định hướng mục tiêu/kết quả (goal-based /
  outcome-based jobs), có sự khác biệt rõ rệt theo từng nhóm ngành
  (Banking, Fintech, Retail\ldots).
\end{itemize}

\vspace{0.5em}\noindent\textbf{Sự khác biệt giữa các dự án trong nhà trường và doanh nghiệp
- ``Same skills, different
game''}\par

\begin{itemize}
\tightlist
\item
  Đối với nhà trường, mục tiêu là học cách triển khai hệ thống hoặc giải
  quyết bài toán
\item
  Đối với doanh nghiệp, mục tiêu là làm cho khách hàng hài lòng, và tối
  ưu chi phí, thời gian triển khai
\item
  Không có tài liệu hay hướng dẫn rõ ràng về yêu cầu, không rõ các bài
  toán cần giải quyết, cần chủ động tìm hiểu
\item
  Sự khác biệt giữa tutorial và production:

  \begin{itemize}
  \tightlist
  \item
    Tutorial: Bạn là người tạo ra các quy tắc
  \item
    Production: Bạn phải tuân theo các quy tắc do người khác đặt ra
  \end{itemize}
\item
  Chia sẻ kinh nghiệm của speaker khi trải qua các công ty khác nhau, từ
  enterprise đến startup.
\end{itemize}

\vspace{0.5em}\noindent\textbf{Kĩ năng giao tiếp trong môi trường công sở và làm việc với
khách
hàng}\par

\begin{itemize}
\tightlist
\item
  Cần sự tương tác tích cực giữa bạn, các thành viên và khách hàng
\item
  Kĩ năng giao tiếp là 1 trong những kĩ năng quan trọng nhất trong môi
  trường công sở
\item
  Cần phải rõ ràng trong giao tiếp, đặc biệt là đối với công việc, cần
  phải trình bày rõ vấn đề và cách giải quyết
\end{itemize}

\vspace{0.5em}\noindent\textbf{Những kĩ năng và công cụ cần thiết trong thời đại
AI}\par

\begin{itemize}
\tightlist
\item
  Cần phải tạo ra sự khác biệt so với đám đông
\item
  Cần có tính ``Integrity'' - tức là cần làm tốt nhất dự án của mình có
  thể.
\item
  Giới thiệu 5 lợi ích của công việc:

  \begin{itemize}
  \tightlist
  \item
    Lương thưởng
  \item
    Phát triển kỹ năng
  \item
    Mạng lưới quan hệ
  \item
    Kiến thức học được
  \item
    Tiềm năng phát triển \textgreater{} Đừng chỉ quan tâm tìm một công
    việc có mức lương cao, hãy quan tâm thêm 4 yếu tố kia để có một công
    việc tốt và phù hợp với bản thân sau này.
  \end{itemize}
\end{itemize}

\vspace{0.5em}\noindent\textbf{Những Gì Học
Được}\par

\vspace{0.5em}\noindent\textbf{Yêu cầu của thị trường việc làm hiện
nay}\par

\begin{itemize}
\tightlist
\item
  \textbf{Nền tảng cốt lõi}: vẫn là quan trọng trong thời đại AI ngày
  nay
\item
  \textbf{Thị trường việc làm thay đổi từng ngày}: cần phải thay đổi bản
  thân để đáp ứng yêu cầu
\item
  \textbf{Các side-project cần thay đổi}: nhờ có AI, các project cần đạt
  yêu cầu cao hơn về chất lượng, có thể chạy được trong môi trường
  production. \#\#\#\# Khác biệt giữa nhà trường và doanh nghiệp
\item
  Cần tư duy cao hơn so với đồ án trong trường: không chỉ là giải quyết
  vấn đề, mà còn phải tối ưu chi phí, thời gian, và làm cho khách hàng
  hài lòng.
\item
  Cần xây dựng 1 kế hoạch thật kỹ lưỡng: cần có roadmap rõ ràng, phải
  tìm hiểu kĩ về domain, phải tìm hiểu về các tính năng của ứng dụng, và
  phải tuân theo các quy tắc của doanh nghiệp.
\end{itemize}

\vspace{0.5em}\noindent\textbf{Kĩ năng mềm}\par

\begin{itemize}
\tightlist
\item
  Kĩ năng giao tiếp luôn là quan trọng nhất trong môi trường doanh
  nghiệp
\item
  Đôi khi không cần phải quá cứng nhắc đối với công việc, mà cần có sự
  linh hoạt, nhằm cân bằng lợi ích đôi bên.
\end{itemize}

\vspace{0.5em}\noindent\textbf{Những kĩ năng và công cụ cần thiết trong thời đại
AI}\par

\begin{itemize}
\tightlist
\item
  AI chỉ là công cụ, không thay thế được suy nghĩ con người.
\item
  Cần phải có những kĩ năng cốt lõi để có thể nhận biết lỗi sai của AI.
\end{itemize}

\vspace{0.5em}\noindent\textbf{Một số hình ảnh khi tham gia sự
kiện}\par

\begin{figure}
\centering
\pandocbounded{\includegraphics[keepaspectratio,alt={Sự kiện AWS}]{aws_report/04_07_1.jpg}}
\caption{Sự kiện AWS}
\end{figure}
